% Forca magnetica sobre uma corrente electrica: 
% circuito rectangular
% Mon Mar  3 21:31:03 WET 2014
% Written by TORDAR 

\documentclass[12pt]{article}  

\pagestyle{empty}

\usepackage{graphicx}
\usepackage{pst-all} 

%Let's go: 

\begin{document}

\null\vskip4cm
\begin{center}
\begin{pspicture}(0,0)(0.5,1)
\psset{unit=8cm}
\psline[linewidth=3pt,linecolor=black]{-}(0.25,0.95)(0.25,0.85)
\psline[linewidth=3pt,linecolor=black]{->}(0.25,0.15)(0.1,0.15)(0.1,0.5)
\psline[linewidth=3pt,linecolor=black]{->}(0.1,0.5)(0.1,0.85)(0.25,0.85)
\psline[linewidth=3pt,linecolor=black]{->}(0.25,0.85)(0.4,0.85)(0.4,0.5)
\psline[linewidth=3pt,linecolor=black]{->}(0.4,0.5)(0.4,0.15)(0.25,0.15)
\psframe[linewidth=3pt,linecolor=blue,framearc=.3](0.5,0.3)
\psframe[linewidth=2pt,linecolor=black,fillstyle=hlines,hatchwidth=0.1pt,hatchsep=3pt](0,0.95)(0.5,1)
\rput{0}(0.05,0.5){\scalebox{1.5}{$I$}}
\rput{0}(0.45,0.5){\scalebox{1.5}{$I$}}
\rput{0}(0.25,0.05){\scalebox{1.5}{$\mathbf{B}$}}
\end{pspicture}
\end{center}

\end{document} 
